\section{Development}\label{sec:development}

GIAANN was developed as a neuralisation of GIA: the symbolic graph machinery of GIA was retained
(substance/action/condition structure, reference delimiters, and concept-level composition), but
execution was shifted from explicit symbolic traversal to sparse neural activation and local synaptic
updates \citep{baxterai_gia,bairesearch2017GIA}. The objective was to preserve GIA's interpretable
semantics while introducing brain-inspired mechanisms for continual learning, scalable memory, and
emergent sequence prediction.

\subsection{GIA History}

The General Intelligence Algorithm (GIA) originated as an explicitly symbolic architecture and was
released publicly through the original OpenGIA SourceForge codebase in 2012
\citep{bairesearch2022OpenGIAsourceforge} (and GitHub in 2017 \citep{bairesearch2017GIA}). That implementation established the core representational
scheme later retained by GIAANN: compositional semantic structures built from substance, action, and
condition roles; explicit reference delimiters for binding and lookup; and graph-based mechanisms for
storing and traversing concept/event relations.

The preliminary Australian provisional patent drafts (1B and 1C) documented this early design as a
general method and computer-system formulation, clarifying the intended scope and operational
workflow of the algorithm \citep{bairesearch2022giaPatentAUprovisional1B,bairesearch2022giaPatentAUprovisional1C}.
Together, the released code and patent drafts provide the historical baseline for GIA as a
semantic-first, interpretable reasoning system, from which the later neuralisation into GIAANN was
developed.

\subsection{From GIA to a Brain-Inspired Neural Substrate}

The early architecture separated a semantic concept network from an episodic instance network,
reflecting GIA's distinction between stable concept knowledge and context-bound events. Over
subsequent revisions, this evolved from a rigid instance ledger toward synaptic event encoding:
episodes became increasingly represented by connection strengths and activation traces over concept
neurons. This was a key transition from symbolic storage to neural memory while keeping the same
conceptual ontology.

Known brain functioning informed each major design choice. Sensory-like feature streams (auditory
and visual) were treated as lower-level inputs converging into a concept integration pathway.
Reference generation (binary/temporal coding) played a thalamus-like indexing role for binding
distributed features into addressable concepts. Later, cortical columnar ideas motivated semantic
concept columns, and sparse competition with selective firing aligned with winner-take-all and
dendritic-style conditional integration \citep{mountcastle1997columnar,london2005dendriticComputation,maass2000winnerTakeAll}.

\subsection{Chronological Architecture Milestones}

The Appendix \Cref{subsec:appendix-giaann-development} GIAANN Development diagrams show the progression below (oldest to newest):

\begin{itemize}
    \item \Cref{fig:dev-2017-initial-core}: initial symbolic core with semantic concepts, episodic instances, and binary reference coding over axonal conduits.
    \item \Cref{fig:dev-2019-event-retrieval}: event retrieval refinement with feedback encoding and rate-coded lookup from partial cues (subject/action/object/place-time subsets).
    \item \Cref{fig:dev-2021-biomodel-mapping}: explicit bio-model mapping across thalamus, cortex, striatum, and limbic/hippocampal roles, including design rationale and predicted serial-access constraints.
    \item \Cref{fig:dev-2021-synaptic-encoding}: transition toward storing episodic information in synapses associated with substance/concept neurons, with a consolidated concept-integrator/reference-generator pathway.
    \item \Cref{fig:dev-2024-columns-routing}: introduction of semantic concept columns and temporary unassociated nodes, with inter-column delimiter features (verbs and prepositions) used as routing constraints.
    \item \Cref{fig:dev-2024-feature-split-gating}: explicit split between inter-column and intra-column features (e.g., qualities), plus syntactic-validity gating so invalid nodes/links remain weak or temporary.
\end{itemize}

\subsection{Resulting GIAANN Design}

Across these stages, GIAANN preserved GIA's semantic compositionality while increasingly grounding
computation in biologically inspired principles: modular columns, sparse competitive activation,
temporally indexed references, and distributed synaptic memory. The result is not a purely symbolic
parser and not a standard backpropagation ANN, but a semantic-first, brain-inspired architecture for
interpretable continual learning and sequence-level reasoning.

\subsection{GIAANN Proto Development}

A full prompt-level natural language specification of GIAANN Proto is provided in Appendix \Cref{subsec:appendix-giaann-proto-spec} (\texttt{dev/GIAANNproto1.nlc}). All debugging prompts are omitted (OP is provided). Manually programmed items have prefix \texttt{MANUAL}.

The major public version of the prototype is \texttt{1}, and the minor public versions are alphabetised (e.g. \texttt{1k}). The large language model used to assist programming the natural language code is provided for each minor public version. There was a significant period of time between the release of OpenAI o1 preview \citep{openai2024o1preview} and GPT-5-Codex \citep{openai2025gpt5codex, openai2025gpt51codex} where all programming had to be performed manually due to the complexity of both the algorithm and the codebase. The majority of programming could be performed in natural language with assistance from GPT-5.2-Codex+ \citep{openai2025gpt52codex, openai2026gpt53codex, openai2026gpt54codex}.